%$Id: config_nux.tex 145 2005-03-25 08:26:35Z myk $

\bghdr{images/fond-nux}

%\begin{center}
%\includegraphics{images/logo_Linux}
%\end{center}

\subsection{Configuration sous Linux}

\subsubsection{Configuration IP}
Quelle que soit ta distribution, il faut que tu décides d'un nom de
machine --- ton pseudo en général. Tu as également besoin de
conna\^itre les informations suivantes :
\begin{itemize}
  \item ton adresse IP %notée dorénavant \server{129.104.AAA.BBB} ;
  \item l'IP de ta passerelle (\emph{gateway})
  \item ton masque de sous-réseau
  \item ton \emph{broadcast}
\end{itemize}
%\begin{description}
%  \item [ton adresse IP] notée dorénavant \server{129.104.AAA.BBB} ;
%  \item [l'IP de ta passerelle (\emph{gateway})] notée \server{129.104.GGG.CCC} ;
%  \item [ton masque de sous-réseau] \server{255.255.FFF.DDD} ;
%  \item [ton broadcast] qui est de la forme \server{129.104.GGG.EEE} avec $EEE = 255 - DDD$. \\
%        \textcolor{red}{(attention, $GGG = AAA$ n'est valable que dans les nouveaux caserts, le PEM et le BEM. Pour Foch, Fayolle %et Maunoury ce n'est donc pas nécessairement vrai que $GGG =
%        AAA$.)}.
%\end{description}
Toutes les informations nécessaires se trouvent page
\pageref{calcul_ip}.

Bien sûr, l'ensemble des manipulations doit se faire en tant que \emph{root}.


\paragraph{\distrib{Sous Gentoo}}

Tu dois d'abord configurer ta Gentoo pour utiliser le réseau de l'X, donc configure le réseau en statique avec ton adresse IP calculée précédemment, puis
configure la résolution DNS. Tu dois avoir ceci dans le fichier \file{/etc/resolv.conf} :
\noindent \cmdline{
search eleves.polytechnique.fr polytechnique.fr\\
nameserver 129.104.201.53\\
nameserver 129.104.201.52}

\label{gentoo_mirror} Pour pouvoir utiliser \app{emerge} à travers le \emph{proxy} de l'école, il faut définir les variables d'environnement
ci-dessous dans le fichier \file{/etc/make.conf} :
\cmdline{http\_proxy=http://kuzh.polytechnique.fr:8080\\
ftp\_proxy=http://kuzh.polytechnique.fr:8080\\
no\_proxy=.eleves.polytechnique.fr\\
GENTOO\_MIRRORS="ftp://miroir.eleves.polytechnique.fr/gentoo http://gentoo.osuosl.org/"\\
SYNC="rsync://rsync/gentoo-portage"}

Tu peux évidemment ajouter d'autres miroirs (séparés par des espaces) dans ta liste mais
\urllink{ftp://miroir} étant interne, il sera toujours beaucoup plus rapide que les autres. On te
conseille ici \urllink{http://gentoo.osuosl.org/} qui est un miroir très complet mais très lent.

\paragraph{\distrib{Ubuntu/Kubuntu}}
\label{Ubuntu:IP}

Lors de l'installation d'une nouvelle Ubuntu, le programme d'installation t'a normalement demand\'e de rentrer toutes ces informations de mani\`ere interactive. Si tu n'as pas fait �ce moment l\`a, tu peux les modifier comme ci-dessous avec ton \'editeur de texte pr\'ef\'er\'e (le tout avec les droits administrateurs \'evidemment!).

\begin{itemize}
\item Le fichier \emph{/etc/hostname} contient ton nom de machine. Il doit contenir uniquement:

\cmdline{tonPseudo.eleves.polytechnique.fr}

\item Le fichier \emph{/etc/resolv.conf} d\'ecrit comment associer le nom d'une machine \`a une adresse IP.
Il doit contenir:

\cmdline{
domain eleves.polytechnique.fr\\
search eleves.polytechnique.fr polytechnique.fr\\
nameserver 129.104.201.53\\
nameserver 129.104.201.52\\
}

\item Le fichier \emph{/etc/network/interfaces} contient entre autres ton IP,
ton sous-r\'eseau et la passerelle pour en sortir.
Ce fichier doit ressembler (avec \'eventuellement une config wifi \`a la suite...,
voir la page \ref{Wifi:linux}) \`a : 

\cmdline{
\# The loopback network interface\\
auto lo\\
iface lo inet loopback\\
\\
\# This is a list of hotpluggable network interfaces.\\
\# They will be activated automatically by the hotplug subsystem.\\
mapping hotplug\\
        script grep\\
        map eth0\\
\\
\# The primary network interface\\
iface eth0 inet static\\
        address   129.104.AAA.BBB\\
        netmask   255.255.255.DDD\\
        network   129.104.AAA.DDD\\
        broadcast 129.104.AAA.EEE\\
        gateway   129.104.AAA.CCC\\
}

\end{itemize}

Ensuite il faut redemarrer ta configuration r\'eseau

\cmdline{\$ sudo /etc/init.d/networking restart}

Voila ta configuration r\'eseau est termin\'ee! Tu peux la tester en pinguant frankiz

\cmdline{\$ ping frankiz}

Tu devrais voir:

\cmdline{
PING frankiz.eleves.polytechnique.fr (129.104.201.51) 56(84) bytes of data.\\
64 bytes from Frankiz.eleves.polytechnique.fr (129.104.201.51) : icmp\_seq=1 ttl=62 time=0.570 ms\\
}

{\bf Configuration du gestionnaire de paquets :}
Il faut d\'esormais configurer le gestionnaire de paquets pour qu'il utilise les miroirs du BR
et non les miroirs \`a l'ext\'erieur du campus qui sont plus lents.

Le fichier \emph{/etc/apt/sources.list} liste les miroirs utilis\'es par le gestionnaire de paquets.
Il faut commenter la premi\'ere ligne (qui correspond au cd d'installation)
ainsi que toutes les lignes non comment\'ees du fichier
(qui correspondent aux miroirs ext\'erieurs au campus) de la fa�n suivante:

\cmdline{deb cdrom:[...]/ breezy main restricted}

devient

\cmdline{\#deb cdrom:[...]/ breezy main restricted}

Il faut ensuite ajouter les lignes suivantes, qui correspondent aux miroirs du BR, au \textbf{d\'ebut} du fichier:

\cmdline{
deb ftp://miroir/ubuntu [version] main restricted universe multiverse\\
deb ftp://miroir/ubuntu [version]-backports main restricted universe multiverse\\
deb ftp://miroir/ubuntu [version]-updates main restricted universe multiverse\\
deb ftp://miroir/ubuntu [version]-security main restricted universe multiverse\\
}

o\`u \textbf{[version]} correspond \`a la version d'ubuntu install\`ee. La version actuelle est \textbf{dapper} et la pr\'ec\'edente est \textbf{breezy}.

On finit par v\'erifier que tout fonctionne en mettant \`a jour la liste des paquets disponibles:

\cmdline{\$ sudo apt-get update}

S'il n'y a pas de message d'erreur c'est que tout fonctionne nickel.


\subsubsection{Configuration antivirus}
\footnotesize{(elle est drôle celle-là hein ?)}

\subsubsection{Configuration du pare-feu}

La solution la plus simple pour se faire un \emph{firewall} sous linux est d'utiliser les \emph{iptables}. Pour ceci la première étape est
d'installer le paquet \app{iptables} pour ta distribution. Pour savoir comment configurer ton \emph{firewall} pour le réseau de l'X, consulte le Wikix.

\subsubsection{Configuration navigateur web}

\flimage{images/nux_firefox_icon}{0.12}{l} Le BR te conseille d'utiliser \app{Firefox} ou
\app{Konqueror} (le navigateur fourni par défaut avec KDE). Dans tous les cas, la seule
configuration à effectuer est celle du \emph{proxy}. Pour \app{Firefox}, il suffit d'aller dans
\menu{Edit}, \menu{Preferences} et dans l'onglet \menu{Avancé}, puis de cliquer sur l'onglet
\menu{Réseau}, et enfin \menu{Paramètres de connexion} ; ensuite coche la case \menu{Adresse de
configuration automatique du proxy}, et rentre \urllink{http://frankiz/proxy.pac}.

Sous \app{Konqueror}, cela se trouve dans le menu \menu{Configuration}, \menu{Configurer Konqueror},
dans l'onglet \menu{Serveur mandataire}. \emph{Attention}: si tu ne configures pas le proxy dans Konqueror,
les logiciels KDE (KGet, Adept, ...) n'utiliserons pas le proxy !!

\imagepos{images/nux_proxy_firefox}{0.35}{Configuration du proxy sous Firefox}{ht}



\subsubsection{Configuration mail}

\flimage{images/nux_kmail_icon}{0.12}{l} Le client mail le plus
utilisé est \app{Kmail}, mais il en existe bien sûr d'autres comme
\app{Thunderbird}. La configuration est semblable, quel que soit le
client que tu utilises.

Va dans \menu{Configuration}, \menu{Configurer Kmail}. Choisis la
rubrique \menu{Comptes}. Commence par créer un nouveau compte dans
l'onglet \menu{Réception des messages} en cliquant sur le bouton
\menu{Ajouter\ldots} et choisis le type POP3.

\imagepos{images/nux_config_kmail_pop}{0.45}{Configuration de la réception des messages sous Kmail}{pht}

Utilise les paramètres suivants pour configurer l'onglet
\menu{Général} :
\begin{description}
  \item[Nom] le nom du compte, par exemple : Mails Poly
  \item[Utilisateur] rentre le login \server{poly} que t'a fourni la DSI à ton arrivée sur le plateau
  \item[Mot de passe] et là le mot de passe \server{poly}
  \item[Serveur] \server{poly.polytechnique.fr}
  \item[Port] 995
\end{description}

Ensuite, va dans l'onglet \menu{Extras} et coche la case
\menu{Utiliser SSL pour sécuriser les téléchargements}.

\imagepos{images/nux_config_kmail_smtp}{0.45}{Configuration de l'envoi des messages sous Kmail}{pht}

Maintenant, dans l'onglet \menu{Envoi des messages} clique sur le
bouton \menu{Ajouter\ldots}. Utilise les paramètres suivants pour le
configurer :
\begin{description}
  \item[Nom] le m\^eme nom de compte que précédemment
  \item[Serveur] \server{poly.polytechnique.fr}
  \item[Port] 25
\end{description}
Sinon, laisse toutes les cases décochées.

\par Tu peux aussi configurer l'accès à \app{l'annuaire LDAP} de l'école, sorte de carnet d'adresses en ligne qui contient les adresses mail de tout le monde sur le campus. Pour ce faire, commence par ouvrir \menu{Outils}, \menu{Carnet d'adresses}, puis va dans \menu{Configuration}, \menu{Configurer kAdressBook}, \menu{Consultation LDAP}. Clique ensuite sur \menu{Ajouter un hôte}, et configure comme suit:
\begin{description}
  \item[Hôte] \server{ldap.eleves.polytechnique.fr}
  \item[Port] 389
  \item[Version de LDAP] 3
  \item[DN] dc=ldap,dc=eleves,dc=polytechnique,dc=fr
  \item[Sécurité] Non
  \item[Identification] Anonyme
\end{description}
Une fois revenu dans \menu{Configuration LDAP}, coche la case \server{ldap.eleves.polytechnique.fr}. Tu as maintenant accès au LDAP lors de la
rédaction de messages, avec tout au plus un redémarrage de \app{Kmail}. \imagepos{images/nux_config_ldap}{0.55}{Configuration de l'annuaire LDAP sous
Kmail}{pht}

\subsubsection{Configuration news}
\flimage{images/nux_knode_icon}{0.12}{l} Le client news le plus
utilisé est \app{Knode}. Parmi les autres clients news, citons
\app{Thunderbird} ou \app{slrn}. Ici aussi, la configuration est
presque indépendante du logiciel choisi.

Sous \app{Knode}, c'est dans le menu \menu{Configuration}, puis \menu{Configurer Knode}. Va dans la rubrique \menu{Comptes, Forums de discussion} et
crée un compte en cliquant sur \menu{Ajouter\ldots}.

\imagepos{images/nux_config_knode}{0.45}{Configuration de Knode}{ht}

Remplis l'onglet \menu{Serveur} avec les informations suivantes :
\begin{description}
  \item[Nom] ce que tu veux pour décrire ce compte, par exemple 'News Frankiz'
  \item[Serveur] \server{frankiz.polytechnique.fr}
  \item[Port] 119
\end{description}
Ensuite occupe-toi de l'onglet \menu{Identité} :
\begin{description}
  \item[Nom] mets ton pseudo dans ce champ
  \item[Organisation] X, école Polytechnique, comme tu le sens
  \item[Adresse électronique] ton adresse mail, pour que les gens puissent te répondre par mail.
\end{description}

Enfin, pour que \app{Knode} puisse envoyer des mails, il faut aller
dans la rubrique \menu{Comptes}, sous-rubrique \menu{Serveur de
courrier (SMTP)}, et choisir comme serveur d'envoi de mails
\server{poly.polytechnique.fr}, port 25 --- c'est exactement la m\^eme
configuration SMTP que \app{Kmail}.

Si tu veux mettre une signature à la fin des messages que tu
posteras, il te suffit de la mettre dans l'onglet \menu{Identité}.
Sur la plupart des clients la signature est interprétée comme
extérieure au message et n'est en particulier pas incluse dans le
texte cité lorsque tu réponds à un message. Pour définir une
signature à la main, il suffit de mettre \verb*+-- +\ (c'est à dire
-{}-<espace>) sur une ligne, et tout ce qui suivra cette ligne
composera ta signature.

Il ne te reste plus qu'à t'inscrire à des \emph{newsgroups} (reporte-toi à la page \pageref{newsgroups} pour plus d'infos) et à poster !

Pour te connecter aux serveurs de \emph{news} de Polytechnique.org, qui ont un accès sécurisé, avec \app{Knode}, il y a une petite subtilité car il
ne gère pas le SSL. Il faut installer \app{stunnel} qui permet de définir une redirection SSL de port. Dans \file{/etc/stunnel.conf} ou
\file{/etc/stunnel/stunnel.conf} selon ta distribution, mets les lignes suivantes (les trois premières y sont en principe déjà) :

\cmdline{\# location of pid file\\
pid = /etc/stunnel/stunnel.pid\\
\\
\# user to run as\\
setuid = stunnel\\
setgid = stunnel\\
\\
\# Use it for client mode\\
client = yes\\
\\
\# sample service-level configuration\\
\\
{[}nntps{]}\\
accept  = 1119\\
connect = ssl.polytechnique.org:563\\
TIMEOUTclose = 0
}

Il ne te reste plus qu'à lancer \app{stunnel} par :

\cmdline{/etc/init.d/stunnel start}

Et tu peux ainsi lire les news de Polytechnique.org en mettant \server{localhost} comme serveur et
\server{1119} comme port. Il faut aussi que tu coches \menu{Le serveur exige une identification} et
que tu rentres ton nom d'utilisateur à Polytechnique.org et ton mot de passe, que tu peux définir
sur \urllink{http://www.polytechnique.org/acces\_smtp.php}.
